% Caderno 2 — cubo do 2^3 com as respostas nos vértices.
% Dados: `peixoto-2x3` (numeros.md, conjunto 2), médias por ensaio.
% A face direita (temperatura no nível +1) está em azul: é a leitura do
% efeito principal como diferença entre as médias de duas faces opostas.
\documentclass[border=3pt,11pt]{standalone}
% Preâmbulo comum das figuras TikZ das apostilas.
% Carregado por todos os `standalone` deste diretório, para que os cinco
% diagramas não divirjam em fonte, cor ou espessura de traço.
%
% fontenc T1 + inputenc utf8 são obrigatórios: sem eles "ç", "ã" e "é"
% saem trocados ou desaparecem silenciosamente do PDF.
\usepackage[T1]{fontenc}
\usepackage[utf8]{inputenc}
\usepackage[brazil]{babel}
% A ordem importa: `amssymb` antes de `newtxmath`, exatamente como em
% `doeapostila.cls`. Invertida, as duas definem `\Bbbk` e a compilação aborta.
\usepackage{amsmath,amssymb}
\usepackage{newtxtext,newtxmath}   % mesma família serifada da classe doeapostila
\usepackage{tikz}
\usetikzlibrary{arrows.meta,calc,decorations.pathreplacing,positioning,
  shapes.geometric,backgrounds}

% Paleta: preto, cinzas e uma única cor de apoio.
\definecolor{azulDoE}{HTML}{1B4E79}
\definecolor{cinzaDoE}{gray}{0.55}
\definecolor{cinzaClaro}{gray}{0.82}
\definecolor{cinzaFundo}{gray}{0.93}

\tikzset{
  % vértice de plano fatorial, com a resposta escrita dentro
  vertice/.style={circle, draw=black, fill=white, line width=0.5pt,
                  inner sep=0pt, minimum size=8.5mm, font=\scriptsize},
  vertdestaque/.style={circle, draw=azulDoE, fill=azulDoE!10, text=azulDoE,
                  line width=0.9pt, inner sep=0pt, minimum size=8.5mm,
                  font=\scriptsize},
  % ponto de planejamento, sem texto dentro
  ponto/.style={circle, draw=black, fill=black, inner sep=0pt,
                minimum size=2.4mm},
  pontovazio/.style={circle, draw=black, fill=white, line width=0.5pt,
                inner sep=0pt, minimum size=2.6mm},
  pontoazul/.style={circle, draw=azulDoE, fill=azulDoE, inner sep=0pt,
                minimum size=2.8mm},
  axial/.style={rectangle, draw=azulDoE, fill=azulDoE, inner sep=0pt,
                minimum size=2.6mm},
  apagado/.style={circle, draw=cinzaClaro, fill=white, line width=0.5pt,
                inner sep=0pt, minimum size=2.6mm},
  aresta/.style={draw=black, line width=0.5pt},
  arestafundo/.style={draw=cinzaClaro, line width=0.5pt},
  eixo/.style={draw=black, line width=0.5pt, -{Latex[length=1.8mm,width=1.4mm]}},
  seta/.style={draw=black, line width=0.6pt, -{Latex[length=2.2mm,width=1.7mm]}},
  setaazul/.style={draw=azulDoE, line width=0.8pt,
                   -{Latex[length=2.2mm,width=1.7mm]}},
  cota/.style={draw=cinzaDoE, line width=0.5pt,
               {Latex[length=1.6mm]}-{Latex[length=1.6mm]}},
  % Corpos escolhidos para casar com as figuras de matplotlib, que usam
  % 8 pt nos ticks e 9 pt nos rótulos de eixo. Com a base de 11 pt do
  % `standalone` (a mesma da classe doeapostila), \scriptsize = 8 pt e
  % \footnotesize = 9 pt. Abaixo disso a figura deixa de ser legível na
  % página impressa, e é para isso que a folha de prova serve.
  rotulo/.style={font=\footnotesize},
  rotulinho/.style={font=\scriptsize},
  nota/.style={font=\footnotesize, align=left},
  notaazul/.style={font=\footnotesize, align=left, text=azulDoE},
}

\begin{document}
\begin{tikzpicture}[x=1cm, y=1cm]

  % --- geometria do cubo -------------------------------------------------
  % f1 = temperatura  -> direita
  % f2 = catalisador  -> cima
  % f3 = concentração -> fundo (deslocamento diagonal)
  \def\lado{3.4}
  \def\dx{1.5}
  \def\dy{1.25}

  \coordinate (fmm) at (0,0);              % frente: T-, Cat-
  \coordinate (fpm) at (\lado,0);          % frente: T+, Cat-
  \coordinate (fmp) at (0,\lado);          % frente: T-, Cat+
  \coordinate (fpp) at (\lado,\lado);      % frente: T+, Cat+
  \coordinate (tmm) at (\dx,\dy);          % fundo:  T-, Cat-
  \coordinate (tpm) at (\lado+\dx,\dy);
  \coordinate (tmp) at (\dx,\lado+\dy);
  \coordinate (tpp) at (\lado+\dx,\lado+\dy);

  % face de temperatura +1, destacada
  \begin{scope}[on background layer]
    \fill[azulDoE!7] (fpm) -- (tpm) -- (tpp) -- (fpp) -- cycle;
  \end{scope}

  % arestas do fundo, mais claras
  \draw[arestafundo] (tmm) -- (tpm) -- (tpp) -- (tmp) -- cycle;
  \draw[arestafundo] (fmm) -- (tmm);
  % arestas da frente e laterais visíveis
  \draw[aresta] (fmm) -- (fpm) -- (fpp) -- (fmp) -- cycle;
  \draw[aresta] (fpm) -- (tpm);
  \draw[aresta] (fmp) -- (tmp);
  \draw[aresta] (fpp) -- (tpp);

  % --- respostas nos vértices --------------------------------------------
  \node[vertice]      at (fmm) {54,0};
  \node[vertdestaque] at (fpm) {86,5};
  \node[vertice]      at (fmp) {48,0};
  \node[vertdestaque] at (fpp) {63,0};
  \node[vertice]      at (tmm) {63,0};
  \node[vertdestaque] at (tpm) {93,5};
  \node[vertice]      at (tmp) {58,5};
  \node[vertdestaque] at (tpp) {72,0};

  % --- eixos, com os níveis nas pontas das setas -------------------------
  % Postos bem afastados do cubo: rótulo encostado em vértice fica ilegível.
  \coordinate (o) at (-3.15,-1.55);
  \draw[eixo] (o) -- ++(1.25,0)
       node[rotulinho, right, align=left] {temperatura\\ 313 $\to$ 333\,K};
  \draw[eixo] (o) -- ++(0,1.25)
       node[rotulinho, above, align=center] {catalisador\\ A $\to$ B};
  \draw[eixo] (o) -- ++(0.79,0.66)
       node[rotulinho, above right, align=left, xshift=-1mm]
       {concentração\\ 1,0 $\to$ 1,5\,M};

  % --- leitura do efeito principal ---------------------------------------
  \node[nota, anchor=north west] at (-3.15,-2.35)
    {Rendimento médio, em \%, das duas réplicas de cada ensaio.};
  \node[notaazul, anchor=north west] at (-3.15,-2.95)
    {$E_{\mathrm{T}} = \overline{y}_{\text{face }+} - \overline{y}_{\text{face }-}
      = 78{,}750 - 55{,}875 = 22{,}875$};

\end{tikzpicture}
\end{document}
